\begin{center}
    \textbf{AGRADECIMIENTOS}
\end{center}

\vspace{0.25cm}

Agradecemos a nuestro director, el Ph.D. Jaime Andrés Pavlich Mariscal, por la orientación y el acompañamiento que dieron forma a este trabajo. A los profesores del Departamento de Ingeniería de Sistemas, por su formación y por las observaciones que ayudaron a precisar el alcance del proyecto. A nuestras familias y personas cercanas, por el apoyo constante durante estos meses. Al equipo, por el trabajo compartido y por sostener el mismo objetivo cuando el camino se hizo más largo de lo previsto. Y a la Pontificia Universidad Javeriana, por el espacio y los recursos que hicieron posible llevarlo a término.

\vspace{0.5cm}

\noindent\textbf{De parte de Andres Henao Niño}

Le agradezco a mis padres por la compañía, el apoyo y la motivación durante los momentos más difíciles de mi carrera. Por la confianza que me brindaron desde el primer instante y por haberme inculcado los valores que llevaré con orgullo por el resto de mi vida. A mis amigos más cercanos, que conozco desde que comenzó mi carrera, por ser incondicionales y por ser la compañía más honesta y sincera que pude tener durante este periodo de mi vida. Gracias a Joshua, David y Joan por no solo ser mis compañeros de tesis, sino también por haber compartido de su vida personal conmigo y por haberme motivado siempre a exigirme más de lo que yo creía que podía dar, en los momentos más oscuros.

\vspace{0.5cm}

\noindent\textbf{De parte de Joan Emmanuel Orduz Chia}

Quiero agradecer a mis padres por darme el privilegio de poder estar en la universidad de mis sueños, por siempre estar ahí para mi en los momentos más complicados y siempre enseñarme todo en el amor y la empatía. A mis compañeros de tesis, Joshua, Andres y David, por su colaboración en este proceso que no ha sido nada fácil. A mi director Jaime Pavlich, por ser un profesor que estuvo presente en los aprendizajes que más aprecio de la carrera. Al grupo de inducción de la Facultad de Ingeniería, por ayudarme a acceder a tantos momentos especiales y a tantos espacios que me permitieron seguir creciendo en el amor, como persona y profesional; al ahora llamado "Comité de Bienvenida", por darle sentido a mi vida dentro de inducción, por permitirme ver lo lindo que es tenderle la mano a los demás y siempre dar amor sin esperar nada a cambio.

Agradezco también a mis amigos de la universidad, por ser el pilar más importante y quienes le dieron luz a mi camino dentro de la universidad. A David, por las risas que compartimos volviendo a casa en el transmilenio, por todos esos momentos que atesoro tanto al inicio de la carrera donde nos reímos tanto. A Diego, por toda la música que compartimos, por tu compañía en momentos difíciles y por todas las veces que jugamos videojuegos hasta tarde. A Cruz, por ser una de mis mayores inspiraciones por su sabiduría y por su capacidad de acompañar en los momentos más felices y también en los más dificiles. A Isa, por ser la mujer que me motiva a ser un mejor hombre, un mejor amigo, un mejor profesional, un mejor novio y un mejor hijo, gracias por siempre estar ahí, gracias por corregirme cuando me descarrilaba y gracias por darme amor y un pecho en el que descansar cuando lo necesite.

También agradezco a Bogotham Comics por ser un lugar donde encontré reconexión a uno de los hobbies que más feliz me hace. A Nightwing, por enseñarme que la esperanza, la bondad, el amor y la empatía son los pilares de mi vida; y por inspirarme a dar todo mi corazón en todo lo que hago, porque mis heridas no me definen y porque no necesito tener super poderes para poder hacer del mundo un mejor lugar.

Finalmente agradezco a la música, por salvarme, por ayudarme a expresar todo lo que pienso y siento, por darme una voz que puede ser escuchada y por darme la energía para seguir adelante.

\vspace{0.5cm}

\noindent\textbf{De parte de Joshua Prieto Zambrano}

Quiero agradecer, en primer lugar, a mi familia. A mi papá, a mi mamá y a mi hermano, que vieron de cerca todo lo que viví en esta etapa y lo que me costó llegar a entregar este trabajo. A mis gatos, que fueron mi colchón y mi refugio cuando no podía más, y que con solo mirarme ya me acompañaban. A mis amigos de antes de la universidad, Lucas, Juanda, Andrés, David, Migue y muchos más, y a los que conocí en el camino, Toño, Alvis, Kraken, Paz, las bebés de Juanma, Lau C, Tomás R y tantos otros. Son mi segunda familia. Incluyo a Henao, Joan y David, mis compañeros de tesis, que terminaron siendo compañeros de vida.

A mi director, Jaime Pavlich, que fue mucho más que un director: fue un mentor y me hizo enamorarme de esta carrera y de todo lo que tiene para enseñarme. A José Hernando Hurtado y a Juan Pablo Delgado Luna, profesores de mis primeros semestres, que me hicieron entender que yo tenía algo para aportarle a esta carrera y que sí podía vivir de ella. Al grupo de inducción de la Facultad de Ingeniería, mi otra familia. Ahí crecí, me conocí y encontré a mis mejores amigos. A Sistehermosos y a todos los que han sido parte de ese grupo: ahí pude ser yo y conocerme de una forma única, en especial a Cruz, Juseb e Isa, que fueron unos hermanos para mí, y a Vale Salgado por ser mi dupla favorita y mi legado en este grupo. A los Macaquinhos, mi grupo de inducción favorito en donde pude conocer no solo neos increíbles y personas espectaculares, sino también amigos que hasta el día de hoy aprecio y quiero con todo el corazón.

Al Centro Pastoral, por llenar mi espiritualidad y hacer crecer mi amor por la vida y por la humanidad. Ahí le encontré el gusto al servicio, a lo frágil y a lo humano. Gracias en especial a Misión País Colombia, que me volvió más vulnerable y me hizo enamorarme de mi país y de su realidad: los niños de Meseta San Rafael, la belleza de Aposentos, la ternura de Santa Bárbara, los retos de Piñalito y, próximamente, la memoria del Darién. En esos lugares conocí mucho de mí y de la vida.

A Kanye West, por acompañarme con su música en las noches de vacío y en los días de luz. A MF DOOM, por enseñarme lo que es ser original y que puedo ser quien soy sin miedo a lo que piensen los demás. A Hirohiko Araki, por JoJos, que me llenó de emoción y de sonrisas con su historia. A Hidetaka Miyazaki, por los Souls, que me acompañaron durante toda la universidad. Volver a jugarlos era un apoyo en los días malos, y me enseñaron que en el mundo sí existe la luz y que toca luchar aun siendo un sinluz. A Spider-Man, a Gohan del futuro y a Guts, por enseñarme a no rendirme. Lo que ellos son es lo que quiero ser: alguien que sigue luchando por sus propósitos. Al L18 y al B18, por llevarme y traerme a mi casa siempre a salvo, y por acogerme y sostenerme. A San Ignacio de Loyola, por enseñarme tanto con los conceptos que fundó, que hoy me permiten orientar mi vida de una forma más segura, estable y real.

Por último, a mí mismo. Yo soy quien vivió este proceso y sabe por lo que tuvo que pasar, todo lo que aprendió, lo que logró y lo que no. Gracias por ser, por enamorarte y por permanecer en el amor. Eso lo definió todo.

\noindent\textbf{De parte de David Estevan Rodriguez Jurado}

Agradezco a mi familia, por su apoyo incondicional y por enseñarme a ser una mejor persona, gracias a ellos tuve el privilegio de estudiar en la que es para mi la mejor universidad de Colombia y esto es un logro tanto para ellos como para mi. A mis amigos, que me acompañaron en los momentos difíciles y celebraron conmigo los logros alcanzados. A mis compañeros de tesis, Joshua, Henao y Joan, por su colaboración y amistad durante este proceso. A mi director, Jaime Pavlich, por su guía y paciencia. Y a la Pontificia Universidad Javeriana, por brindarme la oportunidad de crecer académica y personalmente.


\vspace{0.5cm}

